\label{chap:pantallas-dualidad}
\index{écran algébrique}
\index{dualité résidu--quotient}

Les entiers \(10\), \(11\), \(24\) et \(26\) apparaissent dans la construction par
des procédures algébriques différentes : longueur d'un code, dimension d'un
module, rang d'un réseau et dimension d'un quotient orthogonal.
L'objectif de cette section est de reconstruire ces procédures et les applications qui
les relient. Une égalité entre entiers ne suffira pas à identifier les
objets qu'ils dénombrent.

Nous appellerons \emph{écran algébrique} un sous-espace ou quotient accompagné
de l'application qui le sélectionne. Cette dénomination n'introduit pas une nouvelle classe
d'espaces : elle abrège le couple formé par l'objet linéaire et son projecteur ou
application quotient. Avec cette convention, les réductions

\[
12\longrightarrow11\longrightarrow10,
\qquad
26\longrightarrow24
\]

sont des énoncés sur des applications définies, non des décompositions numérologiques.
La première chaîne provient d'une représentation de \(A_5\) marquée par le
vecteur dodécaphasique \(K\) ; la seconde, d'une réduction isotrope d'un
réseau lorentzien. Toutes deux s'appuient sur la branche exceptionnelle construite dans
les sections précédentes, mais ne sont pas le même mécanisme.

\section{Le quotient parfait à \(11\) coordonnées}
\label{sec:golay-once}

Soit

\[
C_W=[12,6,6]_3
\]

le code ternaire étendu de Golay. Le poinçonnage d'une coordonnée
produit le code ternaire de Golay

\[
G_{11}=[11,6,5]_3\subset\mathbb F_3^{11}
\]

\cite{MacWilliamsSloane1977}. La longueur onze appartient ici à l'espace
ambiant du code ; elle n'est pas introduite au moyen d'une interprétation physique.

\begin{definition}[Boule ternaire]
Pour \(x\in\mathbb F_3^n\) et \(t\geq0\), la boule de Hamming de rayon \(t\)
centrée en \(x\) est
\[
B_t(x)=\{y\in\mathbb F_3^n:d_H(x,y)\leq t\}.
\]
Son cardinal, indépendant de \(x\), est
\[
V_3(n,t)=\sum_{j=0}^{t}\binom nj2^j.
\]
\end{definition}

\begin{theorem}[Perfection de \(G_{11}\)]
\label{thm:g11-perfecto}
Les boules de Hamming de rayon deux centrées sur les mots de \(G_{11}\)
constituent une partition de \(\mathbb F_3^{11}\). En particulier,
\[
V_3(11,2)=243,
\qquad
|G_{11}|\,V_3(11,2)=3^{11}.
\]
\end{theorem}

\begin{proof}
La distance minimale de \(G_{11}\) est cinq. Si deux boules de rayon deux,
centrées sur des mots distincts \(c,c'\in G_{11}\), avaient un élément
commun \(y\), l'inégalité triangulaire donnerait
\[
d_H(c,c')\leq d_H(c,y)+d_H(y,c')\leq4,
\]
en contradiction avec \(d_H(c,c')\geq5\). Les boules sont donc
disjointes. Leur cardinal est
\[
V_3(11,2)
=\binom{11}{0}+2\binom{11}{1}+2^2\binom{11}{2}
=1+22+220=243=3^5.
\]
Puisque \(G_{11}\) est de dimension six, il contient \(3^6\) mots. En
conséquence,
\[
3^6V_3(11,2)=3^6\cdot3^5=3^{11}
=|\mathbb F_3^{11}|.
\]
La réunion disjointe des boules a le cardinal de l'espace ambiant et
coïncide donc avec lui.
\end{proof}

Le décodage parfait attribue à chaque mot reçu un unique mot
de code à une distance au plus deux. De manière équivalente, l'espace des
syndromes a \(3^{11-6}=243\) éléments. Cette apparition de \(243\) doit
être distinguée de la couche cubique

\[
3\cdot9^2=243
\]

de la complétion \(1000=729+243+27+1\). Une quantité dénombre les syndromes
d'un code ; l'autre dénombre les positions d'une couche de frontière. L'égalité
des cardinaux autorise une comparaison, mais ne fournit pas de bijection
canonique entre les deux réalisations.

Il convient également de séparer la partition métrique

\[
243=1+22+220
\]

de l'identité typée

\[
243=1+22+90+120+10.
\]

Dans la seconde expression, \(90=6\binom62\) et
\(120=8\binom62\) comptent des familles d'incidences dans l'extension
hexade--octade, tandis que \(10\) sera la dimension d'un écran linéaire.
Ce ne sont pas des sous-ensembles de poids \(0,1,2\) d'une même boule de Hamming.
La construction de plans qui sous-tend cette extension s'inscrit dans la
théorie classique des systèmes de Witt \cite{AssmusKey1992}.

\section{\(2\) descendants du code ternaire étendu}
\label{sec:dos-descendientes-cw}

La longueur onze n'est pas le seul rang qui descend de \(C_W\). Le même
code intervient dans le recollement du réseau de racines \(A_2^{12}\), dont le
rang est vingt-quatre. Après fixation des données de chambre et du vecteur de
voisinage construits dans la branche exceptionnelle, le voisin d'indice trois sans
racines est le réseau de Leech \(\Lambda_{24}\). La structure des
dépendances est

\begin{equation}
\label{eq:dos-descendientes-cw}
\begin{gathered}
C_W\xrightarrow{\text{poncturation}}G_{11}\subset\mathbb F_3^{11},\\[2mm]
C_W\xrightarrow{\text{recollement}}N(A_2^{12})
\xrightarrow{\text{voisin d'indice }3}\Lambda_{24}.
\end{gathered}
\end{equation}

La première flèche est un morphisme de codes ; la seconde ligne appartient à la
théorie des réseaux. La source ternaire commune explique la coexistence des
rangs onze et vingt-quatre sans les identifier. En particulier, \(11\) n'est pas
une sous-dimension de \(24\), et \(24\) ne s'obtient pas en complétant le code
poinçonné.

\begin{proposition}[Bifurcation des rangs]
La construction \eqref{eq:dos-descendientes-cw} produit deux invariants de
types distincts :
\[
\operatorname{length}(G_{11})=11,
\qquad
\operatorname{rank}(\Lambda_{24})=24.
\]
Leur provenance commune est le code \(C_W\) ; aucune des deux égalités
n'implique un isomorphisme entre l'espace des coordonnées du code et
l'espace réel engendré par le réseau.
\end{proposition}

\begin{proof}
Le poinçonnage élimine l'une des douze positions et ne modifie pas la dimension
six du code, d'où la longueur onze. Chaque copie de \(A_2\) est de rang
deux ; c'est pourquoi \(A_2^{12}\), ses recollements d'indice fini et ses voisins de
rang plein sont de rang vingt-quatre. Les objets vivent respectivement
dans les catégories des codes sur \(\mathbb F_3\) et des réseaux entiers. Un
isomorphisme entre eux ne serait même pas bien typé sans un foncteur
supplémentaire.
\end{proof}

\section{La représentation icosaédrique de dimension \(12\)}
\label{sec:representacion-a5-doce}

La réduction \(12\to11\to10\) provient d'une seconde réalisation du nombre
douze. Soit \(V_{12}=\mathbb R^{12}\) le module de permutation de \(A_5\) sur
les douze sommets d'un icosaèdre. Notons \(\rho\) cette représentation
et \(\chi_{12}\) son caractère. Pour fixer sans ambiguïté la correspondance
avec les douze positions dodécaphasiques, nous réalisons les sommets comme les classes latérales
à gauche \(A_5/H\), où \(A_5\) agit sur \(\{0,1,2,3,4\}\) et
\[
 H=\langle(0\ 1\ 2\ 3\ 4)\rangle\cong C_5.
\]
La position \(j\) est identifiée à la classe latérale \(r_jH\), pour les
représentants suivants, écrits en notation des images :
\[
\begin{array}{c|c@{\qquad}c|c@{\qquad}c|c}
j&r_j&j&r_j&j&r_j\\ \hline
1&(3,2,4,1,0)&5&(3,2,0,4,1)&9&(3,1,2,4,0)\\
2&(3,1,0,2,4)&6&(4,0,1,2,3)&10&(2,1,4,3,0)\\
3&(0,1,3,4,2)&7&(4,3,2,1,0)&11&(3,4,1,2,0)\\
4&(1,4,2,3,0)&8&(0,2,3,1,4)&12&(4,2,1,3,0)
\end{array}
\]
Ce tableau est la donnée d'orientation qui permet d'appliquer un projecteur de
\(A_5\) au vecteur \(K\) ; sans lui, une norme en coordonnées dépendrait d'une
permutation non spécifiée. Dans l'ordre des classes de conjugaison

\[
1A,\quad2A,\quad3A,\quad5A,\quad5B,
\]

les tailles des classes et le nombre de points fixes sont

\[
\begin{array}{c|rrrrr}
 &1A&2A&3A&5A&5B\\ \hline
|C|&1&15&20&12&12\\
\chi_{12}&12&0&0&2&2.
\end{array}
\]

En effet, les rotations d'ordres deux et trois ne fixent pas de sommets, tandis
que chaque rotation d'ordre cinq fixe les deux sommets de son axe. Écrivons

\[
\phi=\frac{1+\sqrt5}{2},
\qquad
\phi'=\frac{1-\sqrt5}{2},
\qquad
\phi+\phi'=1.
\]

La partie pertinente de la table des caractères de \(A_5\) est

\[
\begin{array}{c|rrrrr}
 &1A&2A&3A&5A&5B\\ \hline
\chi_1&1&1&1&1&1\\
\chi_3&3&-1&0&\phi&\phi'\\
\chi_{3'}&3&-1&0&\phi'&\phi\\
\chi_4&4&0&1&-1&-1\\
\chi_5&5&1&-1&0&0.
\end{array}
\]

\begin{theorem}[Décomposition du module icosaédrique]
\label{thm:descomposicion-a5}
La représentation de permutation sur les douze sommets se décompose sous la forme
\[
V_{12}\cong V_1\oplus V_3\oplus V_{3'}\oplus V_5.
\]
Le sous-espace orthogonal au vecteur uniforme est
\[
V_{11}:=\mathbf1^\perp
\cong V_3\oplus V_{3'}\oplus V_5
\]
et est de dimension onze.
\end{theorem}

\begin{proof}
La multiplicité d'un caractère irréductible \(\chi\) dans \(\chi_{12}\) est
\[
\langle\chi_{12},\chi\rangle
=\frac1{60}\sum_C|C|\chi_{12}(C)\chi(C).
\]
Seuls \(1A,5A,5B\) contribuent. Pour les caractères de dimensions un, trois,
trois prime, quatre et cinq, on obtient respectivement
\[
\frac{12+24+24}{60}=1,
\]
\[
\frac{36+24(\phi+\phi')}{60}=1,
\qquad
\frac{36+24(\phi'+\phi)}{60}=1,
\]
\[
\frac{48-24-24}{60}=0,
\qquad
\frac{60}{60}=1.
\]
Les multiplicités sont donc \((1,1,1,0,1)\). L'unique vecteur fixé par
toutes les permutations est le vecteur uniforme
\(\mathbf1=(1,\ldots,1)\). Son complément orthogonal contient les trois
termes non triviaux et est de dimension \(3+3+5=11\).
\end{proof}

Le projecteur orthogonal sur \(V_{11}\) est explicite :

\[
P_{11}=I_{12}-\frac1{12}J_{12},
\]

où \(J_{12}\) est la matrice dont toutes les entrées valent un. Ainsi, la première
réduction est l'application

\[
\mathbb R^{12}\longrightarrow V_{11},
\qquad
x\longmapsto x-\frac1{12}\Bigl(\sum_{j=1}^{12}x_j\Bigr)\mathbf1.
\]

Cette formule spécifie l'information éliminée : exclusivement la
composante uniforme. L'égalité \(12=1+11\) est une conséquence de la
décomposition du module, non une convention dimensionnelle.

\section[Polarisation du secteur tridimensionnel]{Polarisation du secteur tridimensionnel déterminée par K}
\label{sec:reduccion-once-diez}

Considérons le vecteur dodécaphasique déjà obtenu,

\[
K=(234,543,140,729,659,824,621,58,914,794,146,601),
\qquad
\sum_{j=1}^{12}K_j=6263.
\]

Sa composante de somme nulle est

\[
\widetilde K=P_{11}K
=K-\frac{6263}{12}\mathbf1.
\]

Pour sélectionner l'une des deux représentations tridimensionnelles, nous utilisons la
bijection position--classe latérale fixée dans le tableau précédent. Le projecteur central
correspondant s'écrit

\[
P_{\mathrm{ico},3}=\frac{3}{60}\sum_{g\in A_5}\chi_3(g^{-1})\rho(g).
\]

La somme contient des matrices de permutation et des coefficients dans
\(\mathbb Q(\sqrt5)\) ; par conséquent, \(P_{\mathrm{ico},3}\) est un opérateur exact sur
ce corps, satisfait \(P_{\mathrm{ico},3}^2=P_{\mathrm{ico},3}=P_{\mathrm{ico},3}^{\mathsf T}\) et est de rang trois.

\begin{definition}[Direction polarisante]
La projection icosaédrique du vecteur et la droite qu'elle engendre sont
\[
u_K=P_{\mathrm{ico},3}\widetilde K,
\qquad
\ell_K=\mathbb R u_K.
\]
\end{definition}

\begin{lemma}[Non-dégénérescence de la polarisation]
\label{lem:norma-uk}
Pour le vecteur \(K\) et la correspondance position--classe latérale déclarée,
\[
\|u_K\|^2
=\frac{6638585+2275584\sqrt5}{20}>0.
\]
En particulier, \(\ell_K\) est une droite bien définie.
\end{lemma}

\begin{proof}
En substituant l'expression de \(P_{\mathrm{ico},3}\) dans
\(\|u_K\|^2=\widetilde K^{\mathsf T}P_{\mathrm{ico},3}\widetilde K\), la somme se réduit à
\[
\frac3{60}\sum_{g\in A_5}
\chi_3(g^{-1})\,
\widetilde K^{\mathsf T}\rho(g)\widetilde K.
\]
Chaque produit scalaire est rationnel, puisque \(\rho(g)\) est une matrice de
permutation et \(\widetilde K\in\mathbb Q^{12}\). La séparation des deux
classes d'ordre cinq ne laisse comme seule partie irrationnelle qu'un multiple de
\(\sqrt5\). La somme exacte des soixante termes est
\[
\frac{6638585}{20}+\frac{2275584}{20}\sqrt5.
\]
Les deux coefficients sont positifs et \(\sqrt5>0\) ; la norme ne
s'annule donc pas.
\end{proof}

L'énumération exacte du groupe et des douze classes latérales reconstruit le
projecteur sur \(\mathbb Q(\sqrt5)\), vérifie son idempotence et produit la
norme précédente.

\begin{theorem}[Réduction polarisée \(11\to10\)]
\label{thm:once-diez}
Soit
\[
V_{10}:=(V_3\cap\ell_K^\perp)\oplus V_{3'}\oplus V_5.
\]
Alors
\[
V_{11}=\ell_K\oplus V_{10},
\qquad
\dim V_{10}=10.
\]
Le projecteur sur l'écran de dimension dix est
\[
P_{10}=P_{11}-\frac{u_Ku_K^{\mathsf T}}{\|u_K\|^2}.
\]
\end{theorem}

\begin{proof}
D'après le \cref{lem:norma-uk}, \(u_K\neq0\) et, par définition, \(u_K\in V_3\).
D'où
\[
V_3=\ell_K\oplus(V_3\cap\ell_K^\perp).
\]
En ajoutant les termes orthogonaux \(V_{3'}\) et \(V_5\), le
\cref{thm:descomposicion-a5} donne
\[
V_{11}=\ell_K\oplus
\bigl((V_3\cap\ell_K^\perp)\oplus V_{3'}\oplus V_5\bigr).
\]
La dimension du second terme est \(2+3+5=10\). La formule de \(P_{10}\)
soustrait à \(P_{11}\) le projecteur orthogonal de rang un sur \(\ell_K\), de
sorte que son image est précisément \(V_{10}\).
\end{proof}

\begin{corollary}[Décomposition exacte \(1+3+7\)]
\label{cor:descomposicion-1-3-7}
Soit
\[
W_7(K):=(V_3\cap\ell_K^\perp)\oplus V_5.
\]
Alors
\[
\boxed{
V_{11}=\ell_K\oplus V_{3'}\oplus W_7(K),
\qquad
\dim(\ell_K,V_{3'},W_7)=(1,3,7),}
\]
et
\[
\boxed{
V_{10}=V_{3'}\oplus W_7(K),
\qquad 10=3+7.}
\]
\end{corollary}

\begin{proof}
\(\ell_K\subset V_3\) est une droite, de sorte que
\(\dim(V_3\cap\ell_K^\perp)=2\). Par conséquent,
\(\dim W_7=2+5=7\), et les deux décompositions découlent du
\cref{thm:once-diez}.
\end{proof}

\begin{remark}[Condition supplémentaire pour la lecture \(1+3+6\)]
\label{rem:interfaz-1-3-6}
Les décompositions exactes \(1+3+7\) et \(3+7\) ne sélectionnent pas à elles seules
un temps, un espace étendu tridimensionnel ni un secteur compact de rang
six. Une réalisation physique doit choisir une droite
\(\ell_t\subset W_7(K)\), démontrer
\[
W_7(K)=\ell_t\oplus W_6,
\]
construire un réseau de rang plein
\(\Lambda_6\subset W_6\), stable sous le transport, interpréter
\(W_6/\Lambda_6\) comme le secteur compact et présenter un entrelaceur entre le
lecteur temporel du TRIT, \(W_6\) et ce réseau. Le même transport doit
préserver séparément les trois directions étendues du terme
\(V_{3'}\), sans les mélanger avec \(\ell_t\) ni avec le quotient compact. Ce n'est que
sous ces hypothèses que l'on obtient les décompositions
\[
\boxed{
V_{11}=\ell_K\oplus\ell_t\oplus V_{3'}\oplus W_6,
\qquad 11=1+1+3+6,}
\]
\[
\boxed{
V_{10}=\ell_t\oplus V_{3'}\oplus W_6,
\qquad 10=1+3+6.}
\]
L'absence d'un choix naturel de \(\ell_t\), de stabilité de
\(\Lambda_6\), de rang plein, de séparation des trois directions
étendues ou de l'entrelaceur réfuterait cette réalisation, sans modifier le
corollaire algébrique \(1+3+7\).
\end{remark}

La réduction n'est pas \(A_5\)-équivariante en l'absence de la donnée marquée. Le
sous-module \(V_3\) est irréductible et ne contient pas de droite \(A_5\)-invariante ;
la direction \(\ell_K\) apparaît après l'introduction de \(K\). La séquence

\[
12\longrightarrow11\longrightarrow10
\]

possède ainsi deux natures distinctes : le passage \(12\to11\) est canonique pour
l'action de permutation, tandis que \(11\to10\) est une polarisation
relative au vecteur \(K\) et à la correspondance position--classe latérale déclarée.

\section{Réduction isotrope de rang \(26\to24\)}
\label{sec:reduccion-26-24}

Soit \(U\) le plan hyperbolique entier de base \(e_+,e_-\) et de matrice de Gram

\[
\begin{pmatrix}0&1\\1&0\end{pmatrix}.
\]

C'est un réseau pair, unimodulaire et de signature \((1,1)\). Le réseau

\[
L_{26}:=\Lambda_{24}\oplus U
\]

est pair, unimodulaire et de signature \((25,1)\) ; par la classification des réseaux
indéfinis pairs unimodulaires, il s'identifie à \(II_{25,1}\)
\cite{ConwaySloane1999}.

L'adjectif \emph{lorentzien} est employé au sens algébrique : il indique la
présence d'une direction négative dans la forme bilinéaire. La terminologie
provient de la géométrie de Lorentz et de Minkowski
\cite{Lorentz1904,Minkowski1909} ; elle n'identifie pas à elle seule le réseau à
un espace-temps physique.

\begin{theorem}[Réduction nulle]
\label{thm:reduccion-nula}
Pour le vecteur isotrope primitif \(e_+\in U\), il existe un isomorphisme
canonique, une fois la décomposition \(L_{26}=\Lambda_{24}\oplus U\) fixée,
\[
e_+^\perp/\mathbb Ze_+\cong\Lambda_{24}.
\]
\end{theorem}

\begin{proof}
Tout élément de \(L_{26}\) s'écrit de manière unique sous la forme
\(x=\lambda+ae_++be_-\), avec
\(\lambda\in\Lambda_{24}\) et \(a,b\in\mathbb Z\). Puisque la somme est
orthogonale et \(\langle e_-,e_+\rangle=1\),
\[
\langle x,e_+\rangle=b.
\]
Par conséquent,
\[
e_+^\perp=\Lambda_{24}\oplus\mathbb Ze_+.
\]
Le quotient par \(\mathbb Ze_+\) élimine le second terme et conserve le
premier, ce qui définit l'isomorphisme affirmé.
\end{proof}

La dimension diminue de deux parce que la condition d'orthogonalité élimine
une direction et que le quotient élimine la direction nulle restante. Il ne s'agit pas
d'une projection orthogonale sur un complément positif : une droite
isotrope est contenue dans son propre orthogonal. C'est la raison algébrique
pour laquelle la procédure correcte est \(e_+^\perp/\mathbb Ze_+\).

La matrice du plan \(U\) coïncide formellement avec la matrice qui échange
les deux secteurs du transport bidirectionnel. L'adopter comme forme de Gram
établit une réalisation concrète en réseau de cet opérateur involutif. La réduction
\cref{thm:reduccion-nula} est alors exacte au sein de cette réalisation ;
la coïncidence matricielle, à elle seule, n'identifie pas le réseau à un
espace-temps physique.

Ce cadre ambiant est également celui employé dans la construction de Borcherds : le
secteur de charge centrale vingt-quatre est combiné au plan lorentzien de
rang deux avant la réduction qui conduit à l'algèbre de Lie monstrueuse
\cite{Borcherds1992}. La relation pertinente ici est le quotient de réseau
explicite ; la théorie des algèbres d'opérateurs de sommets suit le corridor
classique décrit précédemment.


\section{Correspondance icosaédrique de McKay et type \texorpdfstring{\(E_8\)}{E8}}
\label{sec:tres-ochos}

Trois occurrences de l'entier huit interviennent près de cette frontière :

\[
8_{\rm oct}=|O_H|,
\qquad
V_8:=V_{3'}\oplus V_5,
\qquad
\operatorname{rank}(E_8)=8.
\]

La première est le cardinal d'une octade du système de Witt ; la deuxième est
un module réel de \(A_5\) ; la troisième est le rang d'un système de racines. Ce ne
sont pas des instances d'un unique objet. Le pont classique pertinent passe par le
groupe icosaédrique binaire :

\[
A_5\longleftarrow2A_5\subset SU(2)
\xrightarrow{\text{McKay}}\widetilde E_8
\]

\cite{McKay1980}. Dans la correspondance de McKay, les sommets du diagramme
affine codent les représentations irréductibles de \(2A_5\), et les arêtes
décrivent la décomposition du produit tensoriel avec la représentation
bidimensionnelle naturelle. Cette construction explique pourquoi l'icosaèdre et
\(E_8\) sont liés ; elle n'identifie pas une octade de Witt au module
\(V_8\), ni l'un de ceux-ci au réseau \(E_8\).

\begin{proposition}[Séparation catégorielle de trois invariants de valeur huit]
Aucune égalité entre les trois invariants précédents ne définit par elle-même
un entrelaceur. Pour en relier deux, il faut respectivement
une application des ensembles d'incidence vers les modules, des modules de \(A_5\) vers les
modules de \(2A_5\), ou des représentations vers des données de réseaux.
\end{proposition}

\begin{proof}
Le cardinal d'un ensemble, la dimension d'un espace vectoriel et le rang
d'un réseau sont des invariants dans des catégories différentes. Un isomorphisme dans
l'une d'elles n'est pas défini pour les objets des autres. La
correspondance de McKay fournit l'un des foncteurs intermédiaires, mais pas
les autres ; c'est pourquoi la coïncidence de l'entier huit est insuffisante.
\end{proof}

\section{Réduction dimensionnelle et contexte de la théorie M}
\label{sec:interfaz-teoria-m}

Les constructions précédentes produisent deux occurrences exactes de onze. La
première est la longueur du code parfait \(G_{11}\) ; la seconde est la
dimension du module actif \(V_{11}=\mathbf1^\perp\). Ces objets ne
s'identifient pas : l'un est défini sur \(\mathbb F_3\) et l'autre sur
\(\mathbb R\). La réduction à dix appartient exclusivement au second et
dépend de la droite marquée \(\ell_K\).

La supergravité en onze dimensions et le réseau de dualités qui a conduit à la
formulation de la théorie M fournissent le contexte physique dans lequel une direction
d'un espace à onze dimensions peut être compactifiée et produire des théories à dix
dimensions \cite{CremmerJuliaScherk1978,HullTownsend1995,Witten1995,Townsend1995}.
La correspondance dimensionnelle suggérée par la construction HMT s'organise
dans le dictionnaire suivant :

\[
\begin{array}{>{\raggedright\arraybackslash}p{0.42\textwidth}
                >{\raggedright\arraybackslash}p{0.46\textwidth}}
\toprule
\textbf{Objeto algebraico} & \textbf{Dato requerido en una realización física}\\
\midrule
\(V_{11}=\ell_K\oplus V_{10}\) &
una dirección compacta y una pantalla de diez dimensiones con un mapa de
campos definido\\
\((m,w,R_0)\mapsto(w,m,R_0^{-1})\) &
sectores de momento y enrollamiento, radio físico y escala de cuerda\\
\(\Lambda_{24}\oplus U\cong II_{25,1}\) &
una interpretación dinámica del sector reticular lorentziano\\
\(C_W\to G_{11}\) y \(C_W\to\Lambda_{24}\) &
un entrelazador que preserve incidencia, acción y espectro\\
\bottomrule
\end{array}
\]

Dans la supergravité à onze dimensions, une réalisation complète doit inclure
au moins le champ métrique \(g_{MN}\), la trois-forme \(C_{MNP}\), le
gravitino \(\psi_M\), son action et ses transformations de supersymétrie. Pour
que l'écran \(V_{11}=\ell_K\oplus V_{10}\) représente une compactification
physique, il faut une application qui envoie ces champs sur des données sur
\(\ell_K\) et \(V_{10}\) et qui préserve la dynamique et la limite de basse
énergie. Aucun de ces champs n'entre comme hypothèse dans les
\cref{thm:g11-perfecto,thm:once-diez,thm:reduccion-nula,thm:dualidad-visible-memoria}.

Par conséquent, l'interface possède deux niveaux bien différenciés. Au
niveau algébrique sont démontrées la perfection du code, la
décomposition \(12=1+3+3'+5\), la polarisation \(11=1+10\), la réduction
isotrope \(26\to24\) et l'involution de rayon. Au niveau physique, ces
structures fournissent un domaine et des équations de compatibilité pour
construire un dictionnaire avec la théorie M. La seconde tâche ajoute des champs et de la
dynamique ; elle ne modifie pas le premier niveau et n'est pas utilisée pour le démontrer.

\Needspace{24\baselineskip}
\section{Diagramme des quotients et des réductions de rang}
\label{sec:sintesis-pantallas}

Le diagramme des rangs peut être résumé sans perdre ses types :

\[
\begin{array}{ccccc}
&&C_W=[12,6,6]_3&&\\[1mm]
&\swarrow&&\searrow&\\[-1mm]
G_{11}\subset\mathbb F_3^{11}&&&&N(A_2^{12})\longrightarrow\Lambda_{24}\\[2mm]
&&V_{12}=\mathbb R^{12}&&\\[-1mm]
&&\downarrow P_{11}&&\\[-1mm]
&&V_{11}=V_3\oplus V_{3'}\oplus V_5&&\\[-1mm]
&&\downarrow P_{10}(K)&&\\[-1mm]
&&V_{10}&&\\[2mm]
&&II_{25,1}=\Lambda_{24}\oplus U&&\\[-1mm]
&&\downarrow\ e_+^\perp/\mathbb Ze_+&&\\[-1mm]
&&\Lambda_{24}.&&
\end{array}
\]

La flèche \(P_{11}\) élimine le mode uniforme ; \(P_{10}(K)\) élimine la
direction sélectionnée par \(K\) ; la réduction nulle élimine une paire
isotrope au moyen de l'orthogonal et du quotient. L'involution
\(\mathsf T\) agit dans un domaine de réseau élargi et relie des rayons
réciproques. Ce sont les mécanismes mathématiques qui soutiennent l'interface
dimensionnelle. Leur formulation explicite permet de distinguer précisément les
théorèmes de rang des interprétations physiques qui peuvent être édifiées
sur eux.
